% =============================================================================
\chapter{Implementation}
\label{ch:implementation}
% =============================================================================

\section{Development Environment}
\label{sec:impl:env}

\begin{table}[H]\centering
  \caption{Development environment inventory.}\label{tab:env}
  \begin{tabularx}{\textwidth}{l X l X}\toprule
    Layer       & Stack                                  & Version & Justification \\\midrule
    OS (train)  & Windows 11 + WSL2 Ubuntu 22.04         & ---     & convenience \\
    OS (cloud)  & Debian (Railway)                       & 12      & LTS, Docker \\
    Python      & CPython                                & 3.9.x   & TF 2.10 cuDNN compatibility \\
    TensorFlow  & TF + Keras                             & 2.10.0  & last MX150-compatible TF \\
    CUDA/cuDNN  & ---                                     & 11.2 / 8.1 & matches TF 2.10 \\
    MediaPipe   & mediapipe-python                       & 0.10.x  & latest stable \\
    OpenCV      & opencv-python                          & 4.11.0  & MP4 decode \\
    FastAPI     & fastapi                                & 0.104+  & async, OpenAPI \\
    Uvicorn     & uvicorn                                & 0.24+   & ASGI server \\
    Pydantic    & pydantic                               & v2      & strict schemas \\
    Node.js     & LTS                                    & 20.x    & Prisma + Vite \\
    React       & react / react-dom                      & 18.x    & concurrent \\
    TypeScript  & typescript                             & 5.x     & types \\
    Vite        & vite                                   & 5.x     & dev/bundle \\
    Tailwind    & tailwindcss                            & 3.x     & utility CSS \\
    Express     & express                                & 4.x     & auth/chat \\
    Prisma      & prisma                                 & 5.x     & ORM/SQLite \\
    Socket.IO   & socket.io                              & 4.x     & chat transport \\
    Expo        & expo                                   & 50.x    & RN tooling \\
    React Native& react-native                           & 0.73    & cross-platform \\
    TFLite (RN) & react-native-tflite                    & 1.x     & on-device infer \\\bottomrule
  \end{tabularx}
\end{table}

\section{Machine-Learning Pipeline}
\label{sec:impl:ml}

\begin{figure}[H]\centering
  \includegraphics[width=0.95\textwidth]{fig_4-1_training_dataflow.png}
  \caption{Training-time data flow from raw clips through \texttt{.npz} cache, augmentation, training, and checkpointing.}
  \label{fig:trainflow}
\end{figure}

Each training run is encapsulated as a Jupyter notebook. The four primary
notebooks are
\texttt{Letters/ASL Letter (English)/Mediapipe\_Training.ipynb},
\texttt{Letters/ArSL Letter (Arabic)/Mediapipe\_Optimized\_Training.ipynb},
\texttt{Words/ASL Word (English)/ASL\_Word\_Training.ipynb}, and
\texttt{Words/ArSL Word (Arabic)/ArSL\_Word\_Training\_v2.ipynb}.

\section{Backend API (FastAPI)}
\label{sec:impl:api}

\subsection{Folder Layout}

\begin{lstlisting}[language=,caption={Backend file layout.},label={lst:apilayout}]
backend/
  app/
    main.py
    config.py
    models/
      loader.py
      letter_predictor.py
      word_predictor.py
      word_predictor_v2.py
      mode_detector.py
      letter_decoder.py
      word_decoder.py
    routes/
      predict.py
      websocket_route.py
      health.py
    schemas/
      prediction_request.py
      prediction_response.py
      websocket_message.py
  requirements.txt
  Dockerfile
  railway.toml
\end{lstlisting}

\subsection{REST Endpoints}

\begin{table}[H]\centering
  \caption{REST endpoints exposed by the FastAPI backend.}\label{tab:rest}
  \begin{tabularx}{\textwidth}{l l X X}\toprule
    Method & Path & Request & Response \\\midrule
    GET  & \texttt{/health}          & ---                                & status JSON \\
    POST & \texttt{/predict/letter}  & \texttt{\{landmarks: float[63], lang\}} & top-5 letters \\
    POST & \texttt{/predict/word}    & \texttt{\{sequence: float[T][F], lang\}} & top-5 words \\
    GET  & \texttt{/info}            & ---                                & model metadata \\\bottomrule
  \end{tabularx}
\end{table}

\subsection{WebSocket Protocol}

\begin{figure}[H]\centering
  \includegraphics[width=0.95\textwidth]{fig_4-2_websocket_sequence.png}
  \caption{Browser $\leftrightarrow$ FastAPI WebSocket sequence.}
  \label{fig:wsseq}
\end{figure}

\subsection{Letter-Stream Decoder}

\begin{figure}[H]\centering
  \includegraphics[width=0.85\textwidth]{fig_4-3_letter_decoder_state.png}
  \caption{Letter-stream decoder state machine.}
  \label{fig:letterdecoder}
\end{figure}

\subsection{Mode Detector}

The mode detector switches between letter and word inference based on
per-frame motion:
\begin{equation}
  m_t = \tfrac{1}{F}\sum_{f=1}^{F} \lvert x_{t,f} - x_{t-1,f} \rvert,
  \label{eq:motion}
\end{equation}
and chooses letter mode when the running mean of $m_t$ falls below a
threshold $\theta_1$ for $K_1$ consecutive frames.

\section{Web Frontend (React)}
\label{sec:impl:web}

\begin{figure}[H]\centering
  \includegraphics[width=0.95\textwidth]{fig_4-4_react_component_tree.png}
  \caption{React component tree for the web client.}
  \label{fig:reacttree}
\end{figure}

\section{Authentication and Chat (Node.js)}
\label{sec:impl:auth}

\begin{figure}[H]\centering
  \includegraphics[width=0.95\textwidth]{fig_4-5_jwt_auth_flow.png}
  \caption{JWT-based authentication flow.}
  \label{fig:jwt}
\end{figure}

The Node.js layer manages user accounts, conversational chat (Socket.IO),
and persistence via Prisma over SQLite (Postgres in production).

\section{Mobile App (React Native + Expo + TFLite)}
\label{sec:impl:mobile}

\begin{figure}[H]\centering
  \includegraphics[width=0.95\textwidth]{fig_4-6_mobile_inference.png}
  \caption{Mobile on-device inference pipeline.}
  \label{fig:mobile}
\end{figure}

\section{Deployment}
\label{sec:impl:deploy}

\begin{figure}[H]\centering
  \includegraphics[width=0.95\textwidth]{fig_4-7_deployment_topology.png}
  \caption{Deployment topology --- Vercel (web), Railway (backend), Expo (mobile).}
  \label{fig:deploy}
\end{figure}

\subsection{Backend Dockerfile}

\begin{lstlisting}[language=bash,caption={Backend Dockerfile (production).},label={lst:dockerfile}]
FROM python:3.9-slim
WORKDIR /app
RUN apt-get update && apt-get install -y libgl1 libglib2.0-0 \
    && rm -rf /var/lib/apt/lists/*
COPY requirements.txt .
RUN pip install --no-cache-dir -r requirements.txt
COPY app/ ./app
COPY model_files/ ./model_files
EXPOSE 8000
CMD ["uvicorn", "app.main:app", "--host", "0.0.0.0", "--port", "8000"]
\end{lstlisting}

\subsection{Environment Variables}

\begin{table}[H]\centering
  \caption{Environment variables used across the three deployable tiers.}\label{tab:envvars}
  \begin{tabularx}{\textwidth}{l l X}\toprule
    Variable & Service & Example \\\midrule
    \texttt{VITE\_API\_URL}  & Web      & \texttt{https://eshara-api.railway.app} \\
    \texttt{VITE\_WS\_URL}   & Web      & \texttt{wss://eshara-api.railway.app/ws/recognize} \\
    \texttt{JWT\_SECRET}     & Node API & 256-bit hex string \\
    \texttt{DATABASE\_URL}   & Node API & \texttt{file:./dev.db} or Postgres URL \\
    \texttt{MODEL\_DIR}      & FastAPI  & \texttt{/app/model\_files} \\
    \texttt{CORS\_ORIGINS}   & FastAPI  & comma-separated origins \\\bottomrule
  \end{tabularx}
\end{table}
